\documentclass[11pt, a4paper]{article}
% ── Shared preamble for all RTOS lab documents ────────────────────────────────
% Usage: % ── Shared preamble for all RTOS lab documents ────────────────────────────────
% Usage: % ── Shared preamble for all RTOS lab documents ────────────────────────────────
% Usage: \input{../preamble}  (from each labNN/ directory)
% ──────────────────────────────────────────────────────────────────────────────

% ── Packages ──────────────────────────────────────────────────────────────────
\usepackage[a4paper, top=2.5cm, bottom=2.5cm, left=2.5cm, right=2.5cm, headheight=28pt]{geometry}
\usepackage[utf8]{inputenc}
\usepackage[T1]{fontenc}
\usepackage{lmodern}
\usepackage{booktabs}
\usepackage{array}
\usepackage{tabularx}
\usepackage{longtable}
\usepackage{multirow}
\usepackage{colortbl}
\usepackage{xcolor}
\usepackage{titlesec}
\usepackage{enumitem}
\usepackage{hyperref}
\usepackage{fancyhdr}
\usepackage{parskip}
\usepackage{listings}
\usepackage[most]{tcolorbox}
\usepackage{fontawesome5}
\usepackage{microtype}
\usepackage{graphicx}
\usepackage{amsmath}

% ── Colors (KMUTNB brand) ──────────────────────────────────────────────────────
\definecolor{kmutnbBlue}{RGB}{0,56,116}
\definecolor{kmutnbGold}{RGB}{186,146,36}
\definecolor{lightGray}{RGB}{245,245,245}
\definecolor{midGray}{RGB}{200,200,200}
\definecolor{codeGray}{RGB}{248,248,248}
\definecolor{codeFrame}{RGB}{220,220,220}
\definecolor{tipteal}{RGB}{0,105,92}
\definecolor{warnAmber}{RGB}{121,85,0}
\definecolor{checkGreen}{RGB}{27,94,32}

% ── Section formatting ─────────────────────────────────────────────────────────
\titleformat{\section}
  {\large\bfseries\color{kmutnbBlue}}
  {\thesection.}{0.5em}{}
  [\color{kmutnbGold}\titlerule]

\titleformat{\subsection}
  {\normalsize\bfseries\color{kmutnbBlue}}
  {\thesubsection.}{0.5em}{}

\titleformat{\subsubsection}
  {\small\bfseries\color{kmutnbBlue}}
  {\thesubsubsection.}{0.5em}{}

\titlespacing*{\section}{0pt}{1.4ex plus .2ex}{0.6ex plus .1ex}
\titlespacing*{\subsection}{0pt}{1.0ex plus .2ex}{0.4ex plus .1ex}

% ── Listings (C and bash) ──────────────────────────────────────────────────────
\lstdefinestyle{cstyle}{
  language        = C,
  basicstyle      = \ttfamily\small,
  keywordstyle    = \bfseries\color{kmutnbBlue},
  commentstyle    = \itshape\color{gray},
  stringstyle     = \color{tipteal},
  numberstyle     = \tiny\color{midGray},
  numbers         = left,
  numbersep       = 6pt,
  stepnumber      = 1,
  backgroundcolor = \color{codeGray},
  frame           = single,
  framerule       = 0.4pt,
  rulecolor       = \color{codeFrame},
  breaklines      = true,
  showstringspaces= false,
  tabsize         = 4,
  xleftmargin     = 1.5em,
  xrightmargin    = 0.5em,
  aboveskip       = 0.8em,
  belowskip       = 0.8em,
}

\lstdefinestyle{bashstyle}{
  language        = bash,
  basicstyle      = \ttfamily\small,
  keywordstyle    = \color{kmutnbBlue},
  commentstyle    = \itshape\color{gray},
  backgroundcolor = \color{black!5},
  frame           = single,
  framerule       = 0.4pt,
  rulecolor       = \color{codeFrame},
  breaklines      = true,
  showstringspaces= false,
  xleftmargin     = 1.5em,
  xrightmargin    = 0.5em,
  aboveskip       = 0.6em,
  belowskip       = 0.6em,
}

\lstdefinestyle{textstyle}{
  basicstyle      = \ttfamily\footnotesize,
  backgroundcolor = \color{black!4},
  frame           = single,
  framerule       = 0.4pt,
  rulecolor       = \color{codeFrame},
  breaklines      = true,
  xleftmargin     = 1.5em,
  xrightmargin    = 0.5em,
  aboveskip       = 0.6em,
  belowskip       = 0.6em,
}

% ── Callout boxes ──────────────────────────────────────────────────────────────
\tcbuselibrary{skins, breakable}

\newtcolorbox{tipbox}[1][]{
  enhanced, breakable,
  colback=tipteal!8, colframe=tipteal, coltitle=white,
  fonttitle=\bfseries\small, title={\faLightbulb\quad Tip#1},
  left=6pt, right=6pt, top=4pt, bottom=4pt,
  boxrule=0.8pt, arc=3pt,
}

\newtcolorbox{warnbox}[1][]{
  enhanced, breakable,
  colback=warnAmber!8, colframe=kmutnbGold, coltitle=white,
  fonttitle=\bfseries\small, title={\faExclamationTriangle\quad Note#1},
  left=6pt, right=6pt, top=4pt, bottom=4pt,
  boxrule=0.8pt, arc=3pt,
}

\newtcolorbox{checkpointbox}[1][]{
  enhanced,
  colback=kmutnbBlue!6, colframe=kmutnbBlue, coltitle=white,
  fonttitle=\bfseries\small, title={\faCheckCircle\quad Checkpoint#1},
  left=6pt, right=6pt, top=4pt, bottom=4pt,
  boxrule=0.8pt, arc=3pt,
}

\newtcolorbox{questionbox}[1][]{
  enhanced, breakable,
  colback=kmutnbGold!10, colframe=kmutnbGold, coltitle=white,
  fonttitle=\bfseries\small, title={\faQuestion\quad Question#1},
  left=6pt, right=6pt, top=4pt, bottom=4pt,
  boxrule=0.8pt, arc=3pt,
}

% ── Checkbox list ──────────────────────────────────────────────────────────────
\newlist{checklist}{itemize}{1}
\setlist[checklist]{label=\faSquare[regular], leftmargin=1.5em, noitemsep}

% ── Misc helpers ───────────────────────────────────────────────────────────────
\newcommand{\file}[1]{\texttt{\small #1}}
\newcommand{\api}[1]{\texttt{\small #1}}

% ── Title block macro ─────────────────────────────────────────────────────────
% Usage: \labtitle{03}{Aperiodic Servers \& Producer--Consumer}{2}{CLO 1 \& CLO 2}
\newcommand{\labtitle}[4]{%
  \begin{center}
    {\color{kmutnbBlue}\rule{\linewidth}{2pt}}\\[6pt]
    {\LARGE\bfseries\color{kmutnbBlue} Laboratory #1}\\[4pt]
    {\large\bfseries #2}\\[2pt]
    {\normalsize Real-Time Operating Systems \enspace\textbullet\enspace
     M.Eng.\ ECE \enspace\textbullet\enspace KMUTNB}\\[8pt]
    {\color{kmutnbGold}\rule{\linewidth}{1pt}}
  \end{center}
  \vspace{0.3cm}
  \begin{tabular}{@{}p{3.8cm} p{10cm}@{}}
    \textbf{Platform}       & NXP FRDM-MCXN236 (ARM Cortex-M33 @ 150\,MHz) \\[2pt]
    \textbf{RTOS}           & FreeRTOS v10.6 \\[2pt]
    \textbf{Toolchain}      & ARM GNU Toolchain (\texttt{arm-none-eabi-gcc}) + Make \\[2pt]
    \textbf{Estimated time} & #3 hours \\[2pt]
    \textbf{CLO mapping}    & #4 \\
  \end{tabular}
  \vspace{0.4cm}
}
  (from each labNN/ directory)
% ──────────────────────────────────────────────────────────────────────────────

% ── Packages ──────────────────────────────────────────────────────────────────
\usepackage[a4paper, top=2.5cm, bottom=2.5cm, left=2.5cm, right=2.5cm, headheight=28pt]{geometry}
\usepackage[utf8]{inputenc}
\usepackage[T1]{fontenc}
\usepackage{lmodern}
\usepackage{booktabs}
\usepackage{array}
\usepackage{tabularx}
\usepackage{longtable}
\usepackage{multirow}
\usepackage{colortbl}
\usepackage{xcolor}
\usepackage{titlesec}
\usepackage{enumitem}
\usepackage{hyperref}
\usepackage{fancyhdr}
\usepackage{parskip}
\usepackage{listings}
\usepackage[most]{tcolorbox}
\usepackage{fontawesome5}
\usepackage{microtype}
\usepackage{graphicx}
\usepackage{amsmath}

% ── Colors (KMUTNB brand) ──────────────────────────────────────────────────────
\definecolor{kmutnbBlue}{RGB}{0,56,116}
\definecolor{kmutnbGold}{RGB}{186,146,36}
\definecolor{lightGray}{RGB}{245,245,245}
\definecolor{midGray}{RGB}{200,200,200}
\definecolor{codeGray}{RGB}{248,248,248}
\definecolor{codeFrame}{RGB}{220,220,220}
\definecolor{tipteal}{RGB}{0,105,92}
\definecolor{warnAmber}{RGB}{121,85,0}
\definecolor{checkGreen}{RGB}{27,94,32}

% ── Section formatting ─────────────────────────────────────────────────────────
\titleformat{\section}
  {\large\bfseries\color{kmutnbBlue}}
  {\thesection.}{0.5em}{}
  [\color{kmutnbGold}\titlerule]

\titleformat{\subsection}
  {\normalsize\bfseries\color{kmutnbBlue}}
  {\thesubsection.}{0.5em}{}

\titleformat{\subsubsection}
  {\small\bfseries\color{kmutnbBlue}}
  {\thesubsubsection.}{0.5em}{}

\titlespacing*{\section}{0pt}{1.4ex plus .2ex}{0.6ex plus .1ex}
\titlespacing*{\subsection}{0pt}{1.0ex plus .2ex}{0.4ex plus .1ex}

% ── Listings (C and bash) ──────────────────────────────────────────────────────
\lstdefinestyle{cstyle}{
  language        = C,
  basicstyle      = \ttfamily\small,
  keywordstyle    = \bfseries\color{kmutnbBlue},
  commentstyle    = \itshape\color{gray},
  stringstyle     = \color{tipteal},
  numberstyle     = \tiny\color{midGray},
  numbers         = left,
  numbersep       = 6pt,
  stepnumber      = 1,
  backgroundcolor = \color{codeGray},
  frame           = single,
  framerule       = 0.4pt,
  rulecolor       = \color{codeFrame},
  breaklines      = true,
  showstringspaces= false,
  tabsize         = 4,
  xleftmargin     = 1.5em,
  xrightmargin    = 0.5em,
  aboveskip       = 0.8em,
  belowskip       = 0.8em,
}

\lstdefinestyle{bashstyle}{
  language        = bash,
  basicstyle      = \ttfamily\small,
  keywordstyle    = \color{kmutnbBlue},
  commentstyle    = \itshape\color{gray},
  backgroundcolor = \color{black!5},
  frame           = single,
  framerule       = 0.4pt,
  rulecolor       = \color{codeFrame},
  breaklines      = true,
  showstringspaces= false,
  xleftmargin     = 1.5em,
  xrightmargin    = 0.5em,
  aboveskip       = 0.6em,
  belowskip       = 0.6em,
}

\lstdefinestyle{textstyle}{
  basicstyle      = \ttfamily\footnotesize,
  backgroundcolor = \color{black!4},
  frame           = single,
  framerule       = 0.4pt,
  rulecolor       = \color{codeFrame},
  breaklines      = true,
  xleftmargin     = 1.5em,
  xrightmargin    = 0.5em,
  aboveskip       = 0.6em,
  belowskip       = 0.6em,
}

% ── Callout boxes ──────────────────────────────────────────────────────────────
\tcbuselibrary{skins, breakable}

\newtcolorbox{tipbox}[1][]{
  enhanced, breakable,
  colback=tipteal!8, colframe=tipteal, coltitle=white,
  fonttitle=\bfseries\small, title={\faLightbulb\quad Tip#1},
  left=6pt, right=6pt, top=4pt, bottom=4pt,
  boxrule=0.8pt, arc=3pt,
}

\newtcolorbox{warnbox}[1][]{
  enhanced, breakable,
  colback=warnAmber!8, colframe=kmutnbGold, coltitle=white,
  fonttitle=\bfseries\small, title={\faExclamationTriangle\quad Note#1},
  left=6pt, right=6pt, top=4pt, bottom=4pt,
  boxrule=0.8pt, arc=3pt,
}

\newtcolorbox{checkpointbox}[1][]{
  enhanced,
  colback=kmutnbBlue!6, colframe=kmutnbBlue, coltitle=white,
  fonttitle=\bfseries\small, title={\faCheckCircle\quad Checkpoint#1},
  left=6pt, right=6pt, top=4pt, bottom=4pt,
  boxrule=0.8pt, arc=3pt,
}

\newtcolorbox{questionbox}[1][]{
  enhanced, breakable,
  colback=kmutnbGold!10, colframe=kmutnbGold, coltitle=white,
  fonttitle=\bfseries\small, title={\faQuestion\quad Question#1},
  left=6pt, right=6pt, top=4pt, bottom=4pt,
  boxrule=0.8pt, arc=3pt,
}

% ── Checkbox list ──────────────────────────────────────────────────────────────
\newlist{checklist}{itemize}{1}
\setlist[checklist]{label=\faSquare[regular], leftmargin=1.5em, noitemsep}

% ── Misc helpers ───────────────────────────────────────────────────────────────
\newcommand{\file}[1]{\texttt{\small #1}}
\newcommand{\api}[1]{\texttt{\small #1}}

% ── Title block macro ─────────────────────────────────────────────────────────
% Usage: \labtitle{03}{Aperiodic Servers \& Producer--Consumer}{2}{CLO 1 \& CLO 2}
\newcommand{\labtitle}[4]{%
  \begin{center}
    {\color{kmutnbBlue}\rule{\linewidth}{2pt}}\\[6pt]
    {\LARGE\bfseries\color{kmutnbBlue} Laboratory #1}\\[4pt]
    {\large\bfseries #2}\\[2pt]
    {\normalsize Real-Time Operating Systems \enspace\textbullet\enspace
     M.Eng.\ ECE \enspace\textbullet\enspace KMUTNB}\\[8pt]
    {\color{kmutnbGold}\rule{\linewidth}{1pt}}
  \end{center}
  \vspace{0.3cm}
  \begin{tabular}{@{}p{3.8cm} p{10cm}@{}}
    \textbf{Platform}       & NXP FRDM-MCXN236 (ARM Cortex-M33 @ 150\,MHz) \\[2pt]
    \textbf{RTOS}           & FreeRTOS v10.6 \\[2pt]
    \textbf{Toolchain}      & ARM GNU Toolchain (\texttt{arm-none-eabi-gcc}) + Make \\[2pt]
    \textbf{Estimated time} & #3 hours \\[2pt]
    \textbf{CLO mapping}    & #4 \\
  \end{tabular}
  \vspace{0.4cm}
}
  (from each labNN/ directory)
% ──────────────────────────────────────────────────────────────────────────────

% ── Packages ──────────────────────────────────────────────────────────────────
\usepackage[a4paper, top=2.5cm, bottom=2.5cm, left=2.5cm, right=2.5cm, headheight=28pt]{geometry}
\usepackage[utf8]{inputenc}
\usepackage[T1]{fontenc}
\usepackage{lmodern}
\usepackage{booktabs}
\usepackage{array}
\usepackage{tabularx}
\usepackage{longtable}
\usepackage{multirow}
\usepackage{colortbl}
\usepackage{xcolor}
\usepackage{titlesec}
\usepackage{enumitem}
\usepackage{hyperref}
\usepackage{fancyhdr}
\usepackage{parskip}
\usepackage{listings}
\usepackage[most]{tcolorbox}
\usepackage{fontawesome5}
\usepackage{microtype}
\usepackage{graphicx}
\usepackage{amsmath}

% ── Colors (KMUTNB brand) ──────────────────────────────────────────────────────
\definecolor{kmutnbBlue}{RGB}{0,56,116}
\definecolor{kmutnbGold}{RGB}{186,146,36}
\definecolor{lightGray}{RGB}{245,245,245}
\definecolor{midGray}{RGB}{200,200,200}
\definecolor{codeGray}{RGB}{248,248,248}
\definecolor{codeFrame}{RGB}{220,220,220}
\definecolor{tipteal}{RGB}{0,105,92}
\definecolor{warnAmber}{RGB}{121,85,0}
\definecolor{checkGreen}{RGB}{27,94,32}

% ── Section formatting ─────────────────────────────────────────────────────────
\titleformat{\section}
  {\large\bfseries\color{kmutnbBlue}}
  {\thesection.}{0.5em}{}
  [\color{kmutnbGold}\titlerule]

\titleformat{\subsection}
  {\normalsize\bfseries\color{kmutnbBlue}}
  {\thesubsection.}{0.5em}{}

\titleformat{\subsubsection}
  {\small\bfseries\color{kmutnbBlue}}
  {\thesubsubsection.}{0.5em}{}

\titlespacing*{\section}{0pt}{1.4ex plus .2ex}{0.6ex plus .1ex}
\titlespacing*{\subsection}{0pt}{1.0ex plus .2ex}{0.4ex plus .1ex}

% ── Listings (C and bash) ──────────────────────────────────────────────────────
\lstdefinestyle{cstyle}{
  language        = C,
  basicstyle      = \ttfamily\small,
  keywordstyle    = \bfseries\color{kmutnbBlue},
  commentstyle    = \itshape\color{gray},
  stringstyle     = \color{tipteal},
  numberstyle     = \tiny\color{midGray},
  numbers         = left,
  numbersep       = 6pt,
  stepnumber      = 1,
  backgroundcolor = \color{codeGray},
  frame           = single,
  framerule       = 0.4pt,
  rulecolor       = \color{codeFrame},
  breaklines      = true,
  showstringspaces= false,
  tabsize         = 4,
  xleftmargin     = 1.5em,
  xrightmargin    = 0.5em,
  aboveskip       = 0.8em,
  belowskip       = 0.8em,
}

\lstdefinestyle{bashstyle}{
  language        = bash,
  basicstyle      = \ttfamily\small,
  keywordstyle    = \color{kmutnbBlue},
  commentstyle    = \itshape\color{gray},
  backgroundcolor = \color{black!5},
  frame           = single,
  framerule       = 0.4pt,
  rulecolor       = \color{codeFrame},
  breaklines      = true,
  showstringspaces= false,
  xleftmargin     = 1.5em,
  xrightmargin    = 0.5em,
  aboveskip       = 0.6em,
  belowskip       = 0.6em,
}

\lstdefinestyle{textstyle}{
  basicstyle      = \ttfamily\footnotesize,
  backgroundcolor = \color{black!4},
  frame           = single,
  framerule       = 0.4pt,
  rulecolor       = \color{codeFrame},
  breaklines      = true,
  xleftmargin     = 1.5em,
  xrightmargin    = 0.5em,
  aboveskip       = 0.6em,
  belowskip       = 0.6em,
}

% ── Callout boxes ──────────────────────────────────────────────────────────────
\tcbuselibrary{skins, breakable}

\newtcolorbox{tipbox}[1][]{
  enhanced, breakable,
  colback=tipteal!8, colframe=tipteal, coltitle=white,
  fonttitle=\bfseries\small, title={\faLightbulb\quad Tip#1},
  left=6pt, right=6pt, top=4pt, bottom=4pt,
  boxrule=0.8pt, arc=3pt,
}

\newtcolorbox{warnbox}[1][]{
  enhanced, breakable,
  colback=warnAmber!8, colframe=kmutnbGold, coltitle=white,
  fonttitle=\bfseries\small, title={\faExclamationTriangle\quad Note#1},
  left=6pt, right=6pt, top=4pt, bottom=4pt,
  boxrule=0.8pt, arc=3pt,
}

\newtcolorbox{checkpointbox}[1][]{
  enhanced,
  colback=kmutnbBlue!6, colframe=kmutnbBlue, coltitle=white,
  fonttitle=\bfseries\small, title={\faCheckCircle\quad Checkpoint#1},
  left=6pt, right=6pt, top=4pt, bottom=4pt,
  boxrule=0.8pt, arc=3pt,
}

\newtcolorbox{questionbox}[1][]{
  enhanced, breakable,
  colback=kmutnbGold!10, colframe=kmutnbGold, coltitle=white,
  fonttitle=\bfseries\small, title={\faQuestion\quad Question#1},
  left=6pt, right=6pt, top=4pt, bottom=4pt,
  boxrule=0.8pt, arc=3pt,
}

% ── Checkbox list ──────────────────────────────────────────────────────────────
\newlist{checklist}{itemize}{1}
\setlist[checklist]{label=\faSquare[regular], leftmargin=1.5em, noitemsep}

% ── Misc helpers ───────────────────────────────────────────────────────────────
\newcommand{\file}[1]{\texttt{\small #1}}
\newcommand{\api}[1]{\texttt{\small #1}}

% ── Title block macro ─────────────────────────────────────────────────────────
% Usage: \labtitle{03}{Aperiodic Servers \& Producer--Consumer}{2}{CLO 1 \& CLO 2}
\newcommand{\labtitle}[4]{%
  \begin{center}
    {\color{kmutnbBlue}\rule{\linewidth}{2pt}}\\[6pt]
    {\LARGE\bfseries\color{kmutnbBlue} Laboratory #1}\\[4pt]
    {\large\bfseries #2}\\[2pt]
    {\normalsize Real-Time Operating Systems \enspace\textbullet\enspace
     M.Eng.\ ECE \enspace\textbullet\enspace KMUTNB}\\[8pt]
    {\color{kmutnbGold}\rule{\linewidth}{1pt}}
  \end{center}
  \vspace{0.3cm}
  \begin{tabular}{@{}p{3.8cm} p{10cm}@{}}
    \textbf{Platform}       & NXP FRDM-MCXN236 (ARM Cortex-M33 @ 150\,MHz) \\[2pt]
    \textbf{RTOS}           & FreeRTOS v10.6 \\[2pt]
    \textbf{Toolchain}      & ARM GNU Toolchain (\texttt{arm-none-eabi-gcc}) + Make \\[2pt]
    \textbf{Estimated time} & #3 hours \\[2pt]
    \textbf{CLO mapping}    & #4 \\
  \end{tabular}
  \vspace{0.4cm}
}


% ── Header / Footer ────────────────────────────────────────────────────────────
\pagestyle{fancy}
\fancyhf{}
\fancyhead[L]{\small\color{kmutnbBlue}\textbf{KMUTNB -- Faculty of Engineering}}
\fancyhead[R]{\small\color{kmutnbBlue}Dept.\ of Electrical and Computer Engineering}
\fancyfoot[L]{\small\color{kmutnbBlue}Real-Time Operating Systems -- M.Eng.\ ECE}
\fancyfoot[C]{\small\thepage}
\fancyfoot[R]{\small\color{kmutnbBlue}Lab 09}
\renewcommand{\headrulewidth}{0.4pt}
\renewcommand{\headrule}{\hbox to\headwidth{%
  \color{kmutnbGold}\leaders\hrule height\headrulewidth\hfill}}
\renewcommand{\footrulewidth}{0.4pt}
\renewcommand{\footrule}{\hbox to\headwidth{%
  \color{midGray}\leaders\hrule height\footrulewidth\hfill}}

\hypersetup{
  colorlinks = true,
  linkcolor  = kmutnbBlue,
  urlcolor   = kmutnbBlue,
  citecolor  = kmutnbBlue,
  pdftitle   = {Lab 09 -- Producer-Consumer IPC with FreeRTOS Queue},
  pdfauthor  = {KMUTNB RTOS Course},
}

% ══════════════════════════════════════════════════════════════════════════════
\begin{document}

\labtitle{09}{Producer-Consumer IPC with FreeRTOS Queue}{4--5}{CLO 5 \& CLO 6}

% ── Objectives ────────────────────────────────────────────────────────────────
\section{Objectives}

By completing this lab you will be able to:

\begin{enumerate}[noitemsep, leftmargin=*]
  \item \textbf{Implement} a high-rate producer task that writes samples into a lock-free
        ring buffer protected by \api{\_\_DMB()} memory barriers.
  \item \textbf{Implement} a consumer task that drains the ring buffer and applies a
        4-tap moving-average FIR filter to each sample.
  \item \textbf{Use} a FreeRTOS queue as a lightweight IPC notification channel between
        the producer and consumer.
  \item \textbf{Measure} end-to-end IPC latency using DWT cycle timestamps.
  \item \textbf{Understand} the design differences between FreeRTOS queue IPC and
        hardware Messaging Unit (MU) IPC on true dual-core MCUs.
\end{enumerate}

\begin{warnbox}
  \textbf{Single-core adaptation.}
  The MCXN236 is a \emph{single-core} Cortex-M33 device (the ``N2'' in the part
  number refers to the MCX N-series family, not core count). This lab implements
  the same producer-consumer pattern and ring-buffer IPC concepts using
  \textbf{FreeRTOS tasks and queues} on a single core. The theoretical
  dual-core/MU content in the Background section shows how the design would extend
  to a true dual-core device (e.g.\ MCXN947).
\end{warnbox}

% ── Prerequisites ─────────────────────────────────────────────────────────────
\section{Prerequisites}

\begin{checklist}
  \item Lab 06 complete (DWT cycle counter working).
  \item ARM GNU Toolchain 14.2 and \texttt{pyocd} working.
  \item NXP MCUXpresso SDK at \texttt{C:/nxp\_sdk/mcuxsdk/mcuxsdk}.
  \item Understanding of memory barriers (\api{\_\_DMB()}, \api{\_\_DSB()}).
\end{checklist}

% ── Background ────────────────────────────────────────────────────────────────
\section{Background}

\subsection{FreeRTOS Queue as IPC Notification}

On a single-core MCU, tasks communicate through FreeRTOS primitives.
This lab uses a \textbf{queue of depth 1} as a notification channel --- the producer
sends a token when new samples are available; the consumer blocks on the queue:

\begin{lstlisting}[style=cstyle]
/* Producer sends notification */
uint32_t token = 1;
xQueueSend(g_notify_queue, &token, 0);   /* non-blocking; drop if full */

/* Consumer blocks for notification */
uint32_t dummy;
xQueueReceive(g_notify_queue, &dummy, portMAX_DELAY);
\end{lstlisting}

\subsection{Ring Buffer with Memory Barriers}

Even on a single core, \api{\_\_DMB()} is the correct idiom for lock-free ring
buffers because it prevents the compiler and CPU pipeline from reordering
the data store before the index update:

\begin{lstlisting}[style=cstyle]
g_ring.data[g_ring.write_idx] = sample;
__DMB();                   /* data visible before index advance */
g_ring.write_idx = next;
\end{lstlisting}

\subsection{Conceptual Dual-Core Reference (MCXN947 / True AMP)}

On a true dual-core device the IPC would use the Messaging Unit (MU) hardware
instead of a FreeRTOS queue, and the ring buffer would live in shared SRAM:

\begin{lstlisting}[style=cstyle]
/* Producer -- on CM33_0 (dual-core device only) */
g_ring.data[g_ring.write_idx] = sample;
__DMB();                        /* all stores visible before index update */
g_ring.write_idx = next;
__DMB();                        /* index visible before MU send */
MU_SendMsg(MUA, 0, next);       /* hardware notification to CM33_1 */

/* Consumer ISR -- on CM33_1 (dual-core device only) */
uint32_t idx = MU_ReceiveMsg(MUB, 0);
__DMB();                        /* MU read complete before we read shared buffer */
process(g_ring.data[idx]);
\end{lstlisting}

The \api{\_\_DMB()} requirements are \emph{stricter} on a multi-core device because
the memory model is weakly ordered across cores.

% ── Project Structure ─────────────────────────────────────────────────────────
\section{Project Structure}

\begin{lstlisting}[style=textstyle]
lab09_freertos_ipc/
  cm33_0/                          <- Main (functional) firmware
    Makefile
    src/
      main_cm33_0.c                <- hardware init, queue + task creation
      adc_producer.c/.h            <- producer task (simulated ADC at 1 kHz)
      fir_consumer.c/.h            <- consumer task (4-tap FIR filter)
      shared_mem.h                 <- RingBuf_t, LatencyTest_t types
      FreeRTOSConfig.h
  cm33_1/                          <- Reference build (conceptual; links cleanly)
    Makefile
    src/
      main_cm33_1.c                <- standalone reference entry point
      fir_consumer.c/.h            <- same consumer logic
      shared_mem.h
\end{lstlisting}

\begin{tipbox}
  Only \file{cm33\_0/} is needed to run the lab on the FRDM-MCXN236.
  \file{cm33\_1/} is a companion reference build to illustrate how the consumer
  would be structured on a separate core.
\end{tipbox}

% ── Part A ─────────────────────────────────────────────────────────────────────
\section{Part A --- Build and Flash}

\subsection{Step 1 --- Build cm33\_0}

\begin{lstlisting}[style=bashstyle]
cd labs/lab09_freertos_ipc/cm33_0
mingw32-make
\end{lstlisting}

Expected output:
\begin{lstlisting}[style=textstyle]
   text    data     bss     dec     hex filename
  28896      20   69852   98768   181d0 lab09_cm33_0.elf
\end{lstlisting}

\subsection{Step 2 --- Flash}

\begin{lstlisting}[style=bashstyle]
mingw32-make flash
\end{lstlisting}

\subsection{Step 3 --- Expected UART Output}

\begin{lstlisting}[style=textstyle]
=== RTOS Lab 09: Dual-Core AMP (CM33_0) ===
[CM33_0] starting scheduler -- free heap: 61440 bytes

[CM33_1] consumed=1000, missed=0, free heap=62440
[CM33_1] consumed=2000, missed=0, free heap=62440
\end{lstlisting}

\begin{checkpointbox}[~A]
  UART showing non-zero consumed count and zero missed samples for at least 10 seconds.
\end{checkpointbox}

% ── Part B ─────────────────────────────────────────────────────────────────────
\section{Part B --- Producer Task}

The producer runs at 1\,kHz using \api{vTaskDelayUntil}. Examine
\file{src/adc\_producer.c}:

\begin{lstlisting}[style=cstyle]
void vADCProducerTask(void *pvParameters)
{
    TickType_t xLastWake = xTaskGetTickCount();
    uint32_t   sample_count = 0;

    for (;;) {
        vTaskDelayUntil(&xLastWake, pdMS_TO_TICKS(1));  /* 1 kHz */

        uint32_t next = (g_ring.write_idx + 1U) % RING_BUF_SIZE;
        if (next == g_ring.read_idx) {
            g_ring.missed++;
            continue;               /* ring full -- drop sample */
        }

        g_ring.data[g_ring.write_idx] = (int16_t)(sample_count % 4096);
        sample_count++;
        __DMB();                    /* data store before index advance */
        g_ring.write_idx = next;

        /* Notify consumer every 16 samples */
        if ((g_ring.write_idx & 0xFU) == 0U) {
            uint32_t token = 1U;
            xQueueSend(g_notify_queue, &token, 0);
        }
    }
}
\end{lstlisting}

\begin{questionbox}[~B1]
  Why does the producer call \api{xQueueSend} only every 16 samples rather than
  every sample? What is the trade-off between notification frequency and latency?
\end{questionbox}

\begin{questionbox}[~B2]
  The producer uses \api{xQueueSend} with timeout 0 (non-blocking). What happens
  to a notification if the consumer has not yet read the previous one from the
  queue? Is this safe for a ring-buffer IPC scheme? Why?
\end{questionbox}

% ── Part C ─────────────────────────────────────────────────────────────────────
\section{Part C --- Consumer Task and FIR Filter}

Examine \file{src/fir\_consumer.c}:

\begin{lstlisting}[style=cstyle]
void vFIRConsumerTask(void *pvParameters)
{
    uint32_t consumed   = 0;
    TickType_t rep_tick = xTaskGetTickCount();

    for (;;) {
        uint32_t dummy;
        xQueueReceive(g_notify_queue, &dummy, pdMS_TO_TICKS(100));

        __DMB();     /* ensure latest write_idx is visible */

        while (g_ring.read_idx != g_ring.write_idx) {
            int16_t raw      = g_ring.data[g_ring.read_idx];
            int32_t filtered = fir_filter(raw);     /* 4-tap moving average */
            g_ring.read_idx  = (g_ring.read_idx + 1U) % RING_BUF_SIZE;
            consumed++;
            (void)filtered;
        }

        /* Report every 1 s */
        if ((xTaskGetTickCount() - rep_tick) >= pdMS_TO_TICKS(1000)) {
            rep_tick = xTaskGetTickCount();
            PRINTF("[CM33_1] consumed=%lu  missed=%lu  heap=%u\r\n",
                   (unsigned long)consumed,
                   (unsigned long)g_ring.missed,
                   (unsigned)xPortGetFreeHeapSize());
        }
    }
}
\end{lstlisting}

\subsection{FIR Filter}

The 4-tap moving-average filter in the consumer:

\begin{lstlisting}[style=cstyle]
static int32_t fir_filter(int16_t new_sample)
{
    static int16_t history[4] = {0};
    static uint8_t idx = 0;

    history[idx] = new_sample;
    idx = (uint8_t)((idx + 1U) % 4U);

    int32_t acc = 0;
    for (uint8_t i = 0; i < 4; i++) acc += history[i];
    return acc / 4;
}
\end{lstlisting}

\begin{checkpointbox}[~C]
  Modify \file{fir\_consumer.c} to \texttt{PRINTF} the running average every 1000
  samples. Show terminal output with filtered values.
\end{checkpointbox}

\begin{questionbox}[~C1]
  The ring buffer uses \texttt{volatile} on \texttt{write\_idx} and
  \texttt{read\_idx}, \emph{and} also uses \api{\_\_DMB()}. What does
  \texttt{volatile} prevent that \api{\_\_DMB()} does not? What does
  \api{\_\_DMB()} prevent that \texttt{volatile} does not?
\end{questionbox}

% ── Part D ─────────────────────────────────────────────────────────────────────
\section{Part D --- IPC Latency Measurement}

Use the DWT cycle counter to measure the latency from when the producer writes the
last sample in a batch to when the consumer first reads it. The \file{shared\_mem.h}
already defines \texttt{LatencyTest\_t} for this:

\begin{lstlisting}[style=cstyle]
/* shared_mem.h */
typedef struct {
    volatile uint32_t send_cycles;   /* producer records DWT->CYCCNT before notify */
    volatile uint32_t recv_cycles;   /* consumer records DWT->CYCCNT on wake-up */
} LatencyTest_t;

extern LatencyTest_t g_latency;
\end{lstlisting}

Add the measurements:

\begin{lstlisting}[style=cstyle]
/* In vADCProducerTask -- just before xQueueSend */
g_latency.send_cycles = DWT->CYCCNT;

/* In vFIRConsumerTask -- first line after xQueueReceive returns */
g_latency.recv_cycles = DWT->CYCCNT;
\end{lstlisting}

Then print the latency in the consumer's periodic report:

\begin{lstlisting}[style=cstyle]
uint32_t lat = g_latency.recv_cycles - g_latency.send_cycles;
PRINTF("[LAT] %lu cycles (~%lu ns @ 150 MHz)\r\n",
       (unsigned long)lat, (unsigned long)(lat * 1000UL / 150UL));
\end{lstlisting}

\begin{checkpointbox}[~D]
  Terminal output showing measured IPC latency in cycles and nanoseconds.
\end{checkpointbox}

\begin{questionbox}[~D1]
  On a true dual-core device, MU hardware latency is typically 50--200 cycles.
  How does your FreeRTOS queue IPC latency compare? What accounts for the
  difference (FreeRTOS context switch overhead, queue ISR path, etc.)?
\end{questionbox}

% ── Deliverables ───────────────────────────────────────────────────────────────
\section{Deliverables}

\renewcommand{\arraystretch}{1.3}
\begin{tabular}{@{} c p{5cm} p{6.5cm} @{}}
  \toprule
  \textbf{\#} & \textbf{Item} & \textbf{Details} \\
  \midrule
  1 & Part A UART & Non-zero consumed, zero missed for 10\,s \\
  2 & Part C UART & Running average printed every 1000 samples \\
  3 & Part D UART & IPC latency in cycles and nanoseconds \\
  4 & Written answers & Questions B1, B2, C1, D1 \\
  5 & \file{shared\_mem.h} listing & Ring buffer + latency struct \\
  \bottomrule
\end{tabular}
\renewcommand{\arraystretch}{1}

% ── Key Concepts ───────────────────────────────────────────────────────────────
\section{Key Concepts Demonstrated}

\renewcommand{\arraystretch}{1.3}
\begin{tabular}{@{} p{6.5cm} p{6cm} @{}}
  \toprule
  \textbf{Concept} & \textbf{Where you saw it} \\
  \midrule
  Lock-free ring buffer with DMB barriers & Parts B \& C \\
  FreeRTOS queue as IPC notification      & Parts B \& C \\
  IPC latency measurement with DWT        & Part D \\
  Producer notification batching          & Part B (Q\,B1) \\
  \texttt{volatile} vs.\ \api{\_\_DMB()} & Part C (Q\,C1) \\
  MU hardware IPC (conceptual reference)  & Background \\
  \bottomrule
\end{tabular}
\renewcommand{\arraystretch}{1}

% ── Reference ─────────────────────────────────────────────────────────────────
\section{Reference}

\renewcommand{\arraystretch}{1.3}
\begin{tabular}{@{} p{6cm} p{6.5cm} @{}}
  \toprule
  \textbf{Resource} & \textbf{Relevant sections} \\
  \midrule
  ARM Cortex-M33 TRM & \S A3.4 Memory ordering, DMB instruction \\
  Barry, \textit{Mastering the FreeRTOS Kernel} &
    Ch.~9 (task notifications), \api{xQueueSend}, \api{xQueueReceive} \\
  NXP MCUXpresso SDK & \file{fsl\_mu.h} (for dual-core reference) \\
  NXP MCXN236 Data Sheet & \S 1 (single-core Cortex-M33 architecture) \\
  \bottomrule
\end{tabular}
\renewcommand{\arraystretch}{1}

\end{document}
